\documentclass[12pt]{article}
\usepackage[utf8]{inputenc}
\usepackage[T1]{fontenc}
\usepackage{hyperref}
\title{Tanimuca-Letuama Language Resources for K-12 Curriculum Development}
\author{Compiled by Claude Sonnet 5 (claude-sonnet-5) for Andr\'as Kornai}
\date{}
\begin{document}
\maketitle

\section{Introduction}
Tanimuca-Letuama (ISO 639-3 tnc; Glottocode tani1257) is an Eastern Tucanoan language of the Apaporis River region of Colombia. Estimates of speaker numbers vary: about 300 in older reference listings and about 700 in the Joshua Project profile for Letuama. Search results turned up little educational material. No school curriculum, primer, dictionary or grammar appeared in them.

\section{Primary Educational Resources}
No primers, textbooks or school programmes for Tanimuca-Letuama appeared in the searches. The language is not known to be taught in schools.

\section{Linguistic and Reference Materials}
\begin{itemize}
\item ISO 639 reference entry for Tanimuca-Retuar\~a (tnc): \url{https://wikipedia.com/wiki/ISO_639:tnc}
\item Wikitongues language page for tnc (EGIDS status, genealogy): \url{https://wikitongues.org/languages/tnc/}
\item Joshua Project, Tanimuca in Colombia (people-group profile): \url{https://joshuaproject.net/people_groups/15257}
\item Joshua Project, Letuama in Colombia: \url{https://joshuaproject.net/people_groups/20306}
\item A Colombian university journal article that lists Tanimuca-Letuama among languages of the Apaporis River (search result only; I did not read its content): \url{https://revistas.unal.edu.co/index.php/formayfuncion/article/download/17281/18122/54745}
\end{itemize}

\section{Religious and Community Resources}
\begin{itemize}
\item ScriptureEarth page for tnc. It lists Scripture portions, read-and-listen text, a songbook and audio of the JESUS Film: \url{https://www.scriptureearth.org/00eng.php?iso=tnc}
\end{itemize}
Joshua Project reports that only Bible portions and audio exist and that no full Bible translation does.

\section{Digital and Audio Resources}
The ScriptureEarth page above carries audio and text. Joshua Project also points to Global Recordings Network, but I did not see a direct GRN URL for this language.

\section{Curriculum Development Context}
Very little published material exists for this language. A K-12 curriculum would have to start with documentation work done with the community: orthography decisions, recording elders, and building a lexicon. Scripture-derived texts could serve as an initial corpus. This should be done with the Tanimuca and Letuama communities and with Colombian indigenous education bodies.
\end{document}